\chapter[Lezione 8 --- 15 ottobre 2026]{Lezione 8: punti isolati, accumulazione all'infinito e limiti di funzioni}
\label{chap:lezione-08}

\section{Punti isolati e accumulazione all'infinito}
\label{sec:lez08-accumulazione-infinito}

La \cref{def:lez07-accumulazione} si completa introducendo
i punti isolati e l'accumulazione nei punti infiniti della
\cref{def:lez07-retta-estesa}.

\subsection{Punti isolati}
\label{sec:lez08-isolati}

\begin{definition}[Punto isolato]
\label{def:lez08-isolato}
Sia \(A\subseteq\R\) non vuoto. Un punto \(x_0\in A\)
si dice \emph{isolato} in \(A\) se esiste \(\delta>0\)
tale che
\[
  D(x_0,\delta)\cap A=\{x_0\}.
\]
\end{definition}

Negando la condizione della \cref{def:lez07-accumulazione},
per ogni \(x_0\in\R\) si ottiene
\[
  x_0\notin A'
  \quad\Longleftrightarrow\quad
  \exists\delta>0:
  \bigl(D(x_0,\delta)\cap A\bigr)
  \setminus\{x_0\}=\varnothing.
\]
Se \(x_0\in A\), questa condizione equivale alla
\cref{def:lez08-isolato}, poiché il centro appartiene
a ogni suo intorno circolare. Si ottiene quindi, per ogni \(x_0\in A\), l'equivalenza
\[
  x_0\text{ è isolato in }A
  \quad\Longleftrightarrow\quad
  x_0\notin A'.
\]
L'ipotesi \(x_0\in A\) è essenziale: si confrontino i punti
\(3\) e \(1\) nell'\cref{ex:lez07-accumulazione-intervallo}:
il primo è isolato in \(A=(1,2)\cup\{3\}\),
mentre il secondo è di accumulazione pur non appartenendo ad \(A\).

\subsection{Punti di accumulazione infiniti}
\label{sec:lez08-accumulazione-estesa}

Gli intorni di \(+\infty\) e \(-\infty\), introdotti nella
\cref{def:lez07-intorni-infinito}, permettono di estendere
la \cref{def:lez07-accumulazione} ai punti infiniti. In questo
caso non occorre escludere il centro dell'intorno, poiché
\(+\infty\) e \(-\infty\) non sono numeri reali.

\begin{definition}[Punti di accumulazione infiniti]
\label{def:lez08-accumulazione-infinito}
Sia \(A\subseteq\R\) non vuoto. Il simbolo \(+\infty\) si dice
\emph{punto di accumulazione} di \(A\) se
\[
  \forall b\in\R,\qquad A\cap(b,+\infty)\ne\varnothing.
\]
Analogamente, \(-\infty\) si dice \emph{punto di accumulazione}
di \(A\) se
\[
  \forall b\in\R,\qquad A\cap(-\infty,b)\ne\varnothing.
\]
\end{definition}

Il collegamento con la \cref{def:limitatezza-superiore-inferiore}
è dato dalla seguente proposizione.

\begin{proposition}[Punti di accumulazione infiniti e limitatezza]
\label{prop:lez08-accumulazione-limitatezza}
Sia \(A\subseteq\R\) non vuoto. Allora:
\begin{enumerate}
  \item \(+\infty\) è un punto di accumulazione di \(A\) se e solo
  se \(A\) non è superiormente limitato;
  \item \(-\infty\) è un punto di accumulazione di \(A\) se e solo
  se \(A\) non è inferiormente limitato.
\end{enumerate}
\end{proposition}

\begin{proof}
\emph{Primo punto.}
Se \(A\) è superiormente limitato, per la \cref{def:maggiorante}
esiste \(\lambda\in\R\) tale che
\[
  a\le\lambda \qquad \forall a\in A.
\]
Scegliendo \(b=\lambda+1\), si ottiene
\[
  A\cap(b,+\infty)=\varnothing.
\]
Per la \cref{def:lez08-accumulazione-infinito},
\(+\infty\) non è quindi di accumulazione.

Viceversa, se \(+\infty\) non è di accumulazione,
per la \cref{def:lez08-accumulazione-infinito}
esiste \(b\in\R\) tale che
\[
  A\cap(b,+\infty)=\varnothing.
\]
Ne segue che \(a\le b\) per ogni \(a\in A\).
Il numero \(b\) è dunque un maggiorante di \(A\),
che risulta superiormente limitato.

\emph{Secondo punto.}
Se \(A\) è inferiormente limitato, per la \cref{def:minorante}
esiste \(\mu\in\R\) tale che
\[
  \mu\le a \qquad \forall a\in A.
\]
Scegliendo \(b=\mu-1\), si ottiene
\[
  A\cap(-\infty,b)=\varnothing.
\]
Per la \cref{def:lez08-accumulazione-infinito},
\(-\infty\) non è quindi di accumulazione.

Viceversa, se \(-\infty\) non è di accumulazione,
per la \cref{def:lez08-accumulazione-infinito}
esiste \(b\in\R\) tale che
\[
  A\cap(-\infty,b)=\varnothing.
\]
Ne segue che \(b\le a\) per ogni \(a\in A\).
Il numero \(b\) è dunque un minorante di \(A\),
che risulta inferiormente limitato.
\end{proof}

\begin{figure}[htbp]
  \centering
  \begin{tikzpicture}[unipd/diagram]
  \draw[->,unipd-muted] (-.8,0) -- (7.5,0) node[right] {\(x\)};
  \draw[very thick,unipd-primary,<-] (-.65,0) -- (2.7,0);
  \draw[unipd-primary,fill=white,thick] (2.7,0) circle (2.4pt);
  \draw[very thick,unipd-accent,->] (4.4,0) -- (7.25,0);
  \draw[unipd-accent,fill=white,thick] (4.4,0) circle (2.4pt);
  \draw[unipd-muted] (2.7,-.1) -- (2.7,.1);
  \draw[unipd-muted] (4.4,-.1) -- (4.4,.1);
  \node[below=6pt] at (2.7,0) {\(\lambda\)};
  \node[below=6pt] at (4.4,0) {\(b\)};
  \node[above=7pt,text=unipd-primary] at (.9,0) {\(A\)};
  \node[above=7pt,text=unipd-accent] at (5.8,0) {\((b,+\infty)\)};
  \node[below=21pt] at (3.55,0) {\(A\cap(b,+\infty)=\varnothing\)};
\end{tikzpicture}

  \caption{Un insieme \(A=(-\infty,\lambda)\) superiormente limitato
  non interseca l'intorno \((b,+\infty)\) di \(+\infty\), con
  \(b>\lambda\). Pertanto \(+\infty\) non è un suo punto di
  accumulazione.}
  \label{fig:lez08-accumulazione-infinito}
\end{figure}

La \cref{prop:lez08-accumulazione-limitatezza} non richiede
l'esistenza di massimi o minimi: ad esempio,
\((-\infty,\lambda)\) non ammette massimo
(\cref{fig:lez08-accumulazione-infinito}).

\section{Definizione generale di limite}
\label{sec:lez08-limiti-definizione}

Gli intorni della \cref{sec:lez07-intorni} e le nozioni di
accumulazione introdotte nelle
\cref{def:lez07-accumulazione,def:lez08-accumulazione-infinito}
permettono ora di definire il limite di una funzione.

\begin{definition}[Limite di una funzione]
\label{def:lez08-limite}
Sia \(f\colon A\longrightarrow\R\), con \(A\subseteq\R\)
non vuoto, e sia \(x_0\in\overline{\R}\) un punto di
accumulazione per \(A\). Si dice che \(f\) \emph{ammette limite}
\(\ell\in\overline{\R}\) per \(x\) che tende a \(x_0\),
e si scrive
\[
  \lim_{x\to x_0}f(x)=\ell,
\]
se, per ogni intorno \(V_\ell\) di \(\ell\), esiste un intorno
\(U_{x_0}\) di \(x_0\) tale che
\[
  f(x)\in V_\ell
  \qquad\forall x\in(U_{x_0}\cap A)\setminus\{x_0\}.
\]
\end{definition}

Non è necessario che \(x_0\in A\); anche quando
\(x_0\in A\), il valore \(f(x_0)\) non influisce sul limite.
Per \(x_0=\pm\infty\), l'esclusione del centro è superflua
poiché \(A\subseteq\R\).

L'ordine dei quantificatori è essenziale: per \emph{ogni}
intorno del valore limite deve esistere un intorno di
\(x_0\) nel quale \emph{tutti} i valori di \(f\) nei
punti ammessi del dominio, distinti da \(x_0\),
appartengano all'intorno assegnato.

\subsection{Limite finito in un punto reale}
\label{sec:lez08-limite-finito}

Nel caso \(x_0,\ell\in\R\), per la \cref{def:lez07-intorno},
si possono scegliere gli intorni circolari
\[
  V_\ell=(\ell-\varepsilon,\ell+\varepsilon),
  \qquad
  U_{x_0}=D(x_0,\delta)=(x_0-\delta,x_0+\delta),
\]
con \(\varepsilon,\delta>0\). La \cref{def:lez08-limite}
assume pertanto la forma
\[
  \forall\varepsilon>0\;\exists\delta_\varepsilon>0:
  \quad
  0<|x-x_0|<\delta_\varepsilon
  \ \Longrightarrow\ |f(x)-\ell|<\varepsilon
  \quad\forall x\in A.
\]
Il raggio \(\delta_\varepsilon\) può dipendere da
\(\varepsilon\), mentre il centro \(x_0\) resta escluso.

La \cref{fig:lez08-limite-finito} rappresenta questa condizione
per una funzione razionale, il cui limite sarà verificato
nell'\cref{ex:lez08-razionale-finito}.

\begin{figure}[htbp]
  \centering
  \begin{tikzpicture}[unipd/diagram,x=1.55cm,y=3.1cm]
  \fill[unipd-soft] (1.5,.5) rectangle (2.5,1);
  \draw[->,unipd-muted] (.35,0) -- (3.72,0) node[right] {\(x\)};
  \draw[->,unipd-muted] (.5,-.04) -- (.5,1.3) node[above] {\(y\)};
  \foreach \xx in {1.5,2,2.5} {
    \draw[dashed,unipd-muted] (\xx,0) -- (\xx,1.09);
  }
  \foreach \yy in {.5,.75,1} {
    \draw[densely dotted,unipd-muted] (.5,\yy) -- (3.4,\yy);
  }
  \draw[very thick,unipd-primary,domain=.72:3.37,samples=90]
    plot (\x,{(2*\x-1)/(\x+2)});
  \fill[unipd-primary] (2,.75) circle (1.9pt);
  \draw[unipd-muted] (1.5,-.025) -- (1.5,.025);
  \draw[unipd-muted] (2,-.025) -- (2,.025);
  \draw[unipd-muted] (2.5,-.025) -- (2.5,.025);
  \node[below=5pt] at (1.5,0) {\(x_0-\delta\)};
  \node[below=5pt] at (2,0) {\(x_0\)};
  \node[below=5pt] at (2.5,0) {\(x_0+\delta\)};
  \node[left=5pt] at (.5,.5) {\(\ell-\varepsilon\)};
  \node[left=5pt] at (.5,.75) {\(\ell\)};
  \node[left=5pt] at (.5,1) {\(\ell+\varepsilon\)};
  \node[above right=2pt,text=unipd-primary] at (2.72,.93) {\(f\)};
\end{tikzpicture}

  \caption{Rappresentazione della condizione di limite finito:
  per \(0<|x-x_0|<\delta\), i valori \(f(x)\) appartengono
  alla fascia \((\ell-\varepsilon,\ell+\varepsilon)\).
  Il grafico illustra \(f(x)=(2x-1)/(x+2)\) in prossimità
  di \(x_0=2\), con \(\ell=3/4\), \(\varepsilon=1/4\) e
  \(\delta=1/2\).}
  \label{fig:lez08-limite-finito}
\end{figure}

\subsection{Altri casi particolari}
\label{sec:lez08-limite-casi}

Poiché \(x_0,\ell\in\overline{\R}\), la
\cref{def:lez08-limite} comprende nove casi, secondo che
ciascuno dei due elementi sia reale, \(+\infty\) o
\(-\infty\). Il caso \(x_0,\ell\in\R\) è stato
esplicitato nella \cref{sec:lez08-limite-finito};
gli altri si ricavano dagli intorni della
\cref{def:lez07-intorni-infinito}.

\begin{exercise}[Forme esplicite della definizione di limite]
\label{ex:lez08-casi-limite}
Scrivere la \cref{def:lez08-limite} per gli altri otto casi
possibili, distinguendo punti di accumulazione finiti e
infiniti e valori limite finiti e infiniti.
\end{exercise}

\begin{proof}[Soluzione dell'esercizio]
Si assume in ogni caso che \(x_0\) sia di accumulazione
per \(A\) e che le implicazioni valgano per ogni \(x\in A\).
Quando \(x_0\in\R\), il punto \(x=x_0\) deve essere escluso.

\emph{Punto \(x_0\in\R\), limite infinito.}
La condizione \(\lim_{x\to x_0}f(x)=+\infty\) equivale a
\[
  \forall M\in\R\;\exists\delta>0:\quad
  0<|x-x_0|<\delta\ \Longrightarrow\ f(x)>M.
\]
Analogamente, per \(\lim_{x\to x_0}f(x)=-\infty\), si richiede
\[
  \forall M\in\R\;\exists\delta>0:\quad
  0<|x-x_0|<\delta\ \Longrightarrow\ f(x)<M.
\]

\emph{Punto \(x_0=+\infty\).}
Per \(\ell\in\R\), si richiede
\[
  \forall\varepsilon>0\;\exists b\in\R:\quad
  x>b\ \Longrightarrow\ |f(x)-\ell|<\varepsilon.
\]
Per i due limiti infiniti si ottengono rispettivamente
\[
  \forall M\in\R\;\exists b\in\R:\quad
  x>b\ \Longrightarrow\ f(x)>M
\]
e
\[
  \forall M\in\R\;\exists b\in\R:\quad
  x>b\ \Longrightarrow\ f(x)<M.
\]

\emph{Punto \(x_0=-\infty\).}
Per \(\ell\in\R\), si richiede
\[
  \forall\varepsilon>0\;\exists b\in\R:\quad
  x<b\ \Longrightarrow\ |f(x)-\ell|<\varepsilon.
\]
Per i due limiti infiniti si ottengono rispettivamente
\[
  \forall M\in\R\;\exists b\in\R:\quad
  x<b\ \Longrightarrow\ f(x)>M
\]
e
\[
  \forall M\in\R\;\exists b\in\R:\quad
  x<b\ \Longrightarrow\ f(x)<M.
\]
In tutti i casi si applica la \cref{def:lez08-limite}
con gli intorni del punto \(x_0\) e del valore limite \(\ell\).
\end{proof}

\section{Esercizi sulla definizione di limite}
\label{sec:lez08-esercizi-limiti}

Nelle verifiche della \cref{def:lez08-limite} occorre
anzitutto controllare il dominio e la condizione
di accumulazione.

\begin{exercise}[Un punto non di accumulazione del dominio]
\label{ex:lez08-radice-meno-infinito}
Stabilire se, nel senso della \cref{def:lez08-limite},
è possibile considerare il limite
\[
  \lim_{x\to-\infty}\frac{2x-1}{\sqrt{x+2}}.
\]
\end{exercise}

\begin{proof}[Soluzione dell'esercizio]
Perché il denominatore sia reale e non nullo, occorre
\[
  x+2>0\quad\Longleftrightarrow\quad x>-2.
\]
Il dominio naturale della funzione è pertanto
\[
  A=(-2,+\infty).
\]
Essendo \(A\) inferiormente limitato, per la
\cref{prop:lez08-accumulazione-limitatezza} il punto
\(-\infty\) non è di accumulazione. Manca quindi
un'ipotesi della \cref{def:lez08-limite}.
\end{proof}

\begin{exercise}[Verifica di un limite finito]
\label{ex:lez08-razionale-finito}
Per la funzione
\[
  f(x)=\frac{2x-1}{x+2},
\]
verificare mediante la \cref{def:lez08-limite} che
\[
  \lim_{x\to2}f(x)=\frac34.
\]
\end{exercise}

\begin{proof}[Soluzione dell'esercizio]
\emph{Dominio e punto di accumulazione.}
La funzione è definita esattamente quando \(x+2\ne0\),
quindi
\[
  A=\R\setminus\{-2\}
   =(-\infty,-2)\cup(-2,+\infty).
\]
Il punto \(2\) è di accumulazione per \(A\). Infatti,
per ogni \(\delta>0\) il numero \(a=2+\delta/2\)
appartiene ad \(A\) e soddisfa
\[
  0<|a-2|=\frac{\delta}{2}<\delta.
\]
Più in generale, ogni punto di \(\overline{\R}\) è di
accumulazione per \(A\): per ogni \(t\in\R\) e
\(\delta>0\), almeno uno dei punti \(t-\delta/2\)
e \(t+\delta/2\) appartiene ad \(A\) ed è distinto
da \(t\). Inoltre, \(A\) è illimitato nei due versi,
quindi \(\pm\infty\) sono di accumulazione per la
\cref{prop:lez08-accumulazione-limitatezza}.

\emph{Verifica delle disuguaglianze.}
Fissato \(\varepsilon>0\), occorre trovare
\(\delta_\varepsilon>0\) tale che
\[
  \left|\frac{2x-1}{x+2}-\frac34\right|<\varepsilon
  \qquad\forall x\in A:\ 0<|x-2|<\delta_\varepsilon.
\]
Equivale al sistema
\begin{equation*}
  \begin{cases}
    \displaystyle\frac34-\varepsilon
      <\frac{2x-1}{x+2}, &\text{(I)},\\[7pt]
    \displaystyle\frac{2x-1}{x+2}
      <\frac34+\varepsilon, &\text{(II)}.
  \end{cases}
\end{equation*}
Si cercherà un raggio \(\delta_\varepsilon\le1\),
in modo che \(|x-2|<\delta_\varepsilon\) implichi
\(x>1\), e quindi \(x+2>0\). In tale intorno il
denominatore è positivo e si possono risolvere le
disuguaglianze senza considerare il ramo \(x<-2\)
del dominio.

Per la disuguaglianza (I), moltiplicando per
\(4(x+2)>0\) si ottiene
\begin{align*}
  (3-4\varepsilon)(x+2)&<4(2x-1)\\
  3x-4\varepsilon x+6-8\varepsilon&<8x-4\\
  10-8\varepsilon&<(5+4\varepsilon)x\\
  x&>\frac{10-8\varepsilon}{5+4\varepsilon}.
\end{align*}
Poiché \(5+4\varepsilon>0\),
\[
  2-\frac{10-8\varepsilon}{5+4\varepsilon}
  =\frac{16\varepsilon}{5+4\varepsilon}>0.
\]
Per la disuguaglianza (II), con lo stesso vincolo \(x+2>0\),
si ha
\begin{align*}
  4(2x-1)&<(3+4\varepsilon)(x+2)\\
  8x-4&<3x+4\varepsilon x+6+8\varepsilon\\
  (5-4\varepsilon)x&<10+8\varepsilon.
\end{align*}
Per risolvere quest'ultima disuguaglianza occorre
considerare il segno di \(5-4\varepsilon\):
\begin{enumerate}
  \item Se \(0<\varepsilon<5/4\), si ottiene
  \[
    x<\frac{10+8\varepsilon}{5-4\varepsilon},
  \]
  e la soglia è maggiore di \(2\), poiché
  \[
    \frac{10+8\varepsilon}{5-4\varepsilon}-2
    =\frac{16\varepsilon}{5-4\varepsilon}>0.
  \]
  \item Se \(\varepsilon=5/4\), la disuguaglianza si
  riduce a \(0<20\), verificata per ogni \(x>1\).
  \item Se \(\varepsilon>5/4\), il coefficiente è
  negativo e il verso si inverte:
  \[
    x>\frac{10+8\varepsilon}{5-4\varepsilon}.
  \]
  La soglia è negativa, quindi la condizione è
  soddisfatta in particolare per ogni \(x>1\).
\end{enumerate}

\emph{Scelta del raggio e conclusione.}
Per scegliere un raggio valido per ogni
\(\varepsilon>0\), si ponga
\[
  \eta:=\min\{\varepsilon,1\}.
\]
Poiché \(0<\eta\le1<5/4\), è sufficiente
soddisfare entrambe le disuguaglianze con \(\eta\)
al posto di \(\varepsilon\). Le soglie appena calcolate
permettono di scegliere
\[
  \delta_\varepsilon:=\min\left\{
    1,\frac{16\eta}{5+4\eta},
      \frac{16\eta}{5-4\eta}
  \right\}>0.
\]
Se \(x\in A\) e \(0<|x-2|<\delta_\varepsilon\), allora
\[
  2-\delta_\varepsilon<x<2+\delta_\varepsilon,
\]
e, per la scelta del raggio,
\[
  x>\frac{10-8\eta}{5+4\eta},
  \qquad
  x<\frac{10+8\eta}{5-4\eta}.
\]
Sono quindi verificate entrambe le disuguaglianze
con \(\eta\) al posto di \(\varepsilon\), e ne segue
\[
  \left|\frac{2x-1}{x+2}-\frac34\right|
  <\eta\le\varepsilon.
\]
Poiché \(\varepsilon>0\) era arbitrario, per la
\cref{def:lez08-limite} si conclude
\[
  \lim_{x\to2}\frac{2x-1}{x+2}=\frac34.
\]
\end{proof}

\begin{exercise}[Verifica di un limite infinito]
\label{ex:lez08-razionale-meno-infinito}
Verificare mediante la \cref{def:lez08-limite} che
\[
  \lim_{x\to-4}\frac{-2}{(x+4)^2}=-\infty.
\]
\end{exercise}

\begin{proof}[Soluzione dell'esercizio]
\emph{Dominio e punto di accumulazione.}
La funzione
\[
  f(x):=\frac{-2}{(x+4)^2}
\]
ha dominio
\[
  A=\R\setminus\{-4\}
   =(-\infty,-4)\cup(-4,+\infty).
\]
Per ogni \(\delta>0\), il punto
\(a=-4+\delta/2\) appartiene ad \(A\) e verifica
\(0<|a+4|<\delta\). Pertanto \(-4\) è di
accumulazione per \(A\), secondo la
\cref{def:lez07-accumulazione}.

\emph{Verifica della definizione.}
Per dimostrare che il limite vale \(-\infty\),
per ogni \(K\in\R\) si deve trovare \(\delta_K>0\)
tale che
\[
  \frac{-2}{(x+4)^2}<K
  \qquad\forall x\in A:\ 0<|x+4|<\delta_K.
\]
Poiché \(x\in A\), si ha \(x\ne-4\) e quindi
\((x+4)^2>0\). Moltiplicando per tale numero,
il verso della disuguaglianza rimane invariato:
\[
  \frac{-2}{(x+4)^2}<K
  \quad\Longleftrightarrow\quad
  -2<K(x+4)^2.
\]

Se \(K\ge0\), quest'ultima disuguaglianza è vera
per ogni \(x\in A\), poiché il membro destro è
non negativo. Basta quindi scegliere, ad esempio,
\(\delta_K=1\).

Se \(K<0\), dividendo per \(K\) si inverte il verso:
\[
  (x+4)^2<\frac{-2}{K}
  \quad\Longleftrightarrow\quad
  |x+4|<\sqrt{\frac{-2}{K}}.
\]
Poiché \(-2/K>0\), è sufficiente scegliere
\[
  \delta_K:=\sqrt{\frac{-2}{K}}>0.
\]
In entrambi i casi esiste il raggio richiesto.
Pertanto, per la \cref{def:lez08-limite}, si conclude
\[
  \lim_{x\to-4}\frac{-2}{(x+4)^2}=-\infty.
\]
\end{proof}
